\subsection{Evaluation Metrics}
\label{sec:method:experiment:metrics}
Four groups of metrics are reported. Accuracy and ranking are reported together and never
separately: on a coarsened EDA graph, CTS-Bench \citep{cts_bench2026} observed the mean
absolute error barely moving while $R^2$ fell below zero. Absolute error alone can look
healthy after the model has lost all discriminative power, which is exactly the failure a
reduction study is most at risk of missing.

\paragraph{Notation.}
Let $\mathcal{D} = \{G_i\}_{i=1}^{N}$ be an evaluation set, $y_i = y(G_i) \in [0, 1]$ the
optimizability target of~\ref{sec:prelim:formalization}, $f_\theta$ the trained encoder and
$R$ the reduction operator, with $R = \operatorname{id}$ for the unreduced baseline.
Residuals are written $r_i = \hat{y}_i - y_i$, and $\bar{y} = N^{-1} \sum_{i} y_i$ is the
target mean over $\mathcal{D}$.

\paragraph{Evaluation protocol.}
Each trained configuration is evaluated in two passes over the test set, differing only in
what the same weights are given as input:
\begin{align}
    \hat{y}_i^{\text{match}} &= f_\theta\!\left( R(G_i) \right),
    \label{eq:method:pass-matched} \\
    \hat{y}_i^{\text{full}} &= f_\theta\!\left( G_i \right).
    \label{eq:method:pass-full}
\end{align}
The matched pass~\eqref{eq:method:pass-matched} sees graphs under the reduction the model
was trained with, and the full-graph pass~\eqref{eq:method:pass-full} sees unreduced graphs.
The matched pass answers RQ3. The pair together answers RQ5, since the difference between
the passes isolates the cost of the structural shift from the cost of the reduction itself.

\subsubsection{Predictive Accuracy}
\label{sec:method:experiment:metrics:accuracy}
The training objective is the smooth L1 loss with transition point $\beta = 0.01$,
\begin{equation}
    \ell_\beta(r) = \begin{cases}
        r^{2} / (2\beta), & |r| < \beta, \\
        |r| - \beta / 2, & \text{otherwise},
    \end{cases}
    \label{eq:method:smoothl1}
\end{equation}
reported on validation and test for comparability with the training curves. Since the
target occupies $[0, 1]$ while $\beta = 0.01$, all but the smallest residuals fall in the
linear branch of~\eqref{eq:method:smoothl1}: outside a band of one percentage point the
objective behaves as a mean absolute error. Accuracy proper is reported as
\begin{align}
    \mathrm{MAE} &= \frac{1}{N} \sum_{i=1}^{N} \left| r_i \right|,
    \label{eq:method:mae} \\
    \mathrm{RMSE} &= \sqrt{\frac{1}{N} \sum_{i=1}^{N} r_i^{2}},
    \label{eq:method:rmse} \\
    R^{2} &= 1 - \frac{\sum_{i=1}^{N} r_i^{2}}
        {\sum_{i=1}^{N} \left( y_i - \bar{y} \right)^{2}}.
    \label{eq:method:r2}
\end{align}
MAE and RMSE are on the label's own scale, so $\mathrm{RMSE} = 0.05$ is five percentage
points of optimizability. By the power mean inequality $\mathrm{MAE} \leq \mathrm{RMSE}$,
with equality exactly when every $\left| r_i \right|$ is identical, so the ratio
$\mathrm{RMSE} / \mathrm{MAE}$ measures how much a few large errors carry the average.
Equation~\eqref{eq:method:r2} divides the same squared error by the target's own spread:
$R^{2} = 0$ is the score of the constant predictor $\hat{y}_i = \bar{y}$, and $R^{2} < 0$ a
model worse than that constant. Because the denominator is $N$ times the sample variance of
$y$ within $\mathcal{D}$, $R^{2}$ is computed for validation and test only, never per
training batch, whose variance is not a meaningful reference for the spread of the target.

\paragraph{Persisted predictions and stratified re-scoring.}
The evaluation sweep writes every per-graph prediction to disk, not only the aggregates, so
every metric above can be recomputed on any subset of the test set from the same forward
pass. This matters because the pooled label is severely skewed: roughly half the test
graphs have $y_i = 0$ exactly and about three quarters have $y_i < 0.005$, so a pooled
$R^2$ is dominated by a mass of near-identical labels and can look respectable while saying
nothing about the graphs a practitioner would ask about. Each accuracy claim is therefore
repeated on three subsets beside the pooled figure: graphs from designs the synthesis
script can still change (two designs sit at a fixed point, with largest observed
optimizability $0.0003$, so they contribute no variance for $R^2$ to explain, only noise
for it to divide by); graphs with $y_i > 0$; and graphs derived from a single upstream
script, which holds constant how much optimization the input received before Orchestrate
saw it. The strata are applied uniformly to every configuration rather than invoked where
they happen to help.

\subsubsection{Ranking Quality}
\label{sec:method:experiment:metrics:ranking}
Spearman rank correlation between predicted and true optimizability across the test set,
that is the Pearson correlation of the midrank vectors,
\begin{equation}
    \rho = \frac{\operatorname{cov}\!\left( \operatorname{rk}(\hat{y}),
        \operatorname{rk}(y) \right)}
        {\sigma_{\operatorname{rk}(\hat{y})} \, \sigma_{\operatorname{rk}(y)}}.
    \label{eq:method:spearman}
\end{equation}
Equation~\eqref{eq:method:spearman} is unchanged by any strictly increasing map applied to
$\hat{y}$, which is why it is reported alongside the error metrics rather than in place of
them: a model whose predictions have collapsed toward $\bar{y}$, the failure mode a
reduction study should expect, holds $\rho$ while $\mathrm{RMSE}$ and $R^2$ degrade.
Ranking is also closer to use. A practitioner deciding where to spend synthesis effort
needs the ordering of circuits, not calibrated values.

The metric ranks \emph{circuits} under one fixed script, not \emph{scripts} for one
circuit. The latter is the script-selection task the motivation invokes and this thesis
does not attempt (\ref{sec:intro:scope:prediction}).

\subsubsection{Efficiency and Resource Metrics}
\label{sec:method:experiment:metrics:efficiency}
Peak device memory, training step and epoch wall-clock time, and inference throughput in
graphs per second. These come from three instruments running under three different batching
regimes, and which instrument produced a number determines what that number may be compared
against, because~\ref{sec:method:experiment:batching} makes budget uniformity a
precondition for reading the efficiency results and that precondition applies to only one
of the three.

\emph{Inference} throughput and peak memory are measured during the evaluation sweep, under
the fixed evaluation node budget of~\ref{sec:method:experiment:batching}, identical across
every configuration. Comparing configurations evaluated at different budgets would be
invalid rather than merely noisy.

\emph{Training step time and step memory} come from a separate benchmark that processes one
graph per batch, deliberately using no budget at all: a node budget packs a different
number of graphs into a batch for every configuration, so per-batch aggregates would
average over incomparable batch compositions. With one graph per batch, each reduced graph
is instead matched to its unreduced counterpart by graph identity, which is what the paired
statistics of~\ref{sec:method:experiment:metrics:stats} require. The price is that these
are per-graph costs, a controlled comparison between configurations rather than an estimate
of wall-clock training cost.

\emph{Epoch wall-clock time} is recorded by the training loop itself under the training
node budget, the smaller of the two. It is comparable across configurations and against
nothing else here.

\paragraph{Memory quantities reported.}
Peak \emph{allocated} memory is the high-water mark of live tensors and the closest measure
of the model's actual demand. Peak \emph{reserved} memory is what the caching allocator has
taken from the driver, is never smaller, and reflects allocation history and fragmentation
as much as demand. Allocated is the primary figure, but reserved sets the out-of-memory
boundary and is recorded alongside it. Both are taken against a \emph{memory floor}
$m_{\mathrm{floor}}$ (parameters, optimizer state and device context, measured after warmup
with no graph resident), so the marginal cost of the graph itself is
\begin{equation}
    \Delta m = m_{\mathrm{peak}} - m_{\mathrm{floor}}.
    \label{eq:method:marginal-memory}
\end{equation}
$\Delta m$ rather than $m_{\mathrm{peak}}$ is the quantity a reduction can move:
$m_{\mathrm{floor}}$ is fixed by the architecture and identical across configurations, so
savings quoted against $m_{\mathrm{peak}}$ would be diluted toward zero by an additive
constant having nothing to do with the reduction.
% TODO: paired savings currently divide raw peak memory, not \eqref{eq:method:marginal-memory}.

\subsubsection{Reduction Quality Metrics}
\label{sec:method:experiment:metrics:reduction}
The node and edge retention ratios $\eta_{\mathcal{V}}$ and $\eta_{\mathcal{E}}$
of~\eqref{eq:prelim:retention} are reported separately rather than collapsed into one
compression figure, together with the offline wall-clock cost of computing the reduction.

The offline cost matters independently of the training saving. Let $c_R$ be the cost of
computing and storing one graph's reduction, and $t_{\operatorname{id}}$ and $t_R$ the
per-graph step cost on the unreduced and the reduced graph. One epoch visits each graph
once, so over $n$ epochs reducing is cheaper than not reducing once
\begin{equation}
    n > n^{\star} = \frac{c_R}{t_{\operatorname{id}} - t_R}.
    \label{eq:method:amortisation}
\end{equation}
Reporting the amortisation threshold $n^{\star}$ rather than a raw offline cost states the
trade-off in the terms that decide it, and makes the failure case explicit:
\eqref{eq:method:amortisation} is undefined for $t_R \geq t_{\operatorname{id}}$, so a
reduction that does not shorten a step never amortises, however cheap it is to compute.
The reduced graph is stored once and reloaded, so $n$ accumulates over every epoch of every
run that reuses the artifact.

\paragraph{Effective receptive field.}
Hypothesis H1 claims that coarsening helps by contracting paths, and that claim needs a
measurement of its own. For a node $v$, let
\begin{equation}
    \mathcal{C}_k(v) = \left\{ u \in \mathcal{V} : u \rightsquigarrow v
        \text{ along a directed path of length at most } k \right\}
    \label{eq:method:cone}
\end{equation}
be its $k$-hop fanin cone, and let $m_u = \mathbf{1}^{\top} x_u$ be the member count of
node $u$, supplied directly by the count-based feature schema
of~\ref{sec:method:architecture} and equal to $1$ on an unreduced graph. The metric is the
member-weighted mean cone size over a node set $\mathcal{S}$ sampled from $\mathcal{V}(G)$,
\begin{equation}
    \bar{c}_k(G) = \frac{1}{\left| \mathcal{S} \right|} \sum_{v \in \mathcal{S}}
        \; \sum_{u \in \mathcal{C}_k(v)} m_u,
    \label{eq:method:receptive}
\end{equation}
evaluated at $k = L$, the number of encoder layers, before and after reduction. The
weighting is what makes the comparison meaningful: counting super-nodes rather than their
members would make $\bar{c}_L$ fall under every coarsening by construction, while weighting
by $m_u$ counts original gates in both states. H1 predicts
$\bar{c}_L(R(G)) > \bar{c}_L(G)$ under summarization and
$\bar{c}_L(R(G)) \leq \bar{c}_L(G)$ under sparsification. Unlike a DAG longest-path metric,
\eqref{eq:method:receptive} is well defined for every method, not only for those that
guarantee an acyclic quotient.
% TODO: not implemented, so H1 is currently asserted rather than evidenced.

\subsubsection{Statistical Treatment}
\label{sec:method:experiment:metrics:stats}
Efficiency comparisons are made pairwise per graph against the full-graph baseline. For
each test graph the saving in an efficiency quantity $q$, either marginal memory $\Delta m$
or step time, is taken relative to that graph's own baseline measurement,
\begin{equation}
    \delta_i = 100 \left( 1 - \frac{q_R(G_i)}{q_{\operatorname{id}}(G_i)} \right),
    \label{eq:method:paired-saving}
\end{equation}
and $\{\delta_i\}$ is summarised by its mean with a $95\%$ percentile bootstrap interval
over $B = 2000$ resamples, together with a two-sided Wilcoxon signed-rank test of
$H_0 \colon \operatorname{med}(\delta) = 0$. Both configurations are measured on the same
graph, so the pairing in~\eqref{eq:method:paired-saving} is by construction, and the
savings are skewed across a corpus whose graph sizes span orders of magnitude, which is why
a rank test replaces a paired $t$-test.

Accuracy comparisons have no equivalent treatment, and this is a real weakness rather than
an omission of detail. Each configuration is trained once, so there is no run-to-run
variance against which to judge whether a difference is real. This bears directly on RQ4,
whose expected effect is small: either the RQ4 pairings are re-run across several seeds, or
the accuracy differences must be reported as point estimates and read with corresponding
caution (\ref{sec:discussion:limitations:reproducibility}).
